% 2011 年北京卷文科第 5 题：四棱锥的三视图。
% 来源 PDF：普通高考/2011/2011北京文.pdf，SHA-256 1ba0def5e75e0a4d7cc4e4e52a11a65ad598b98bc2a8e6c98df8f29fb6b9f4a6。
% 原页：第 1 页（A3 三栏）第 1 栏下部，书页码 649。
% 重绘依据：照原页——正（主）视图为底 4、高 2 的等腰三角形（左侧标 2）；
%   侧（左）视图为同底同高的等腰三角形（原页未标数）；俯视图为边长 4 的正方形，
%   画出自中心连向四个顶点的四条线段（图上呈两条对角线），左边标 4、下边标 4。
%   俯视图与正视图左右对齐（长对正）。不按比例。
\begin{tikzpicture}[>=stealth, semithick, scale=0.9,
  dim/.style={<->, thin}]

  % 正（主）视图：底 4、高 2 的等腰三角形
  \draw (0,0) -- (4,0) -- (2,2) -- cycle;
  % 侧（左）视图：同底同高的等腰三角形
  \draw (6.5,0) -- (10.5,0) -- (8.5,2) -- cycle;

  % 正视图高的标注 2
  \draw[dim] (-0.7,0) -- (-0.7,2);
  \node[left] at (-0.85,1) {$2$};

  % 俯视图：边长 4 的正方形与两条对角线
  \draw (0,-6.0) -- (4,-6.0) -- (4,-2.0) -- (0,-2.0) -- cycle;
  \draw (0,-6.0) -- (4,-2.0);
  \draw (0,-2.0) -- (4,-6.0);

  % 俯视图的纵深 4（左侧）与宽 4（下侧）
  \draw[dim] (-0.7,-6.0) -- (-0.7,-2.0);
  \node[left] at (-0.85,-4.0) {$4$};
  \draw[dim] (0,-6.7) -- (4,-6.7);
  \node[above] at (2,-6.7) {$4$};

  % 视图名称
  \node at (2,-0.85) {正（主）视图};
  \node at (8.5,-0.85) {侧（左）视图};
  \node at (2,-7.45) {俯视图};
\end{tikzpicture}
